\section{Two-walk scales and preliminary comparisons}\label{sec:fluctuations}

We now measure fluctuations in the second coordinate at first hits of
successive height layers. Common true levels provide fresh pair
continuations and, for independent environments, an iid sequence of
symmetric increments. Their truncated variance will determine the
fluctuation scale. We first compare all first-hit positions with this
sequence, then transfer the comparison between two forms of conditioning.
The resulting short-block estimate and endpoint spread bounds will be
used in Sections~\ref{sec:gaussian} and~\ref{sec:clipping}.

On a successful prefix, set
\[
Y_j=(X_{T_j}-X_0)_2.
\]
Choose a deterministic integer median \(b_j\) of \(Y_j\) under \(P_0^+\),
with \(b_0=0\), and write
\[
W_H=\max_{0\le j\le H}|Y_j-b_j|.
\]
All these variables record displacement from the indicated starting site.

\subsection{Common true cuts}

Common spatial regeneration levels for two walks, reached at possibly
different path times, are central to the shared- and independent-environment
comparison of Rassoul-Agha and Sepp\"al\"ainen~\cite[Section~7]{RassoulAghaSeppalainen2009}.
Their invariance principle assumes regeneration-time moments. Here we
establish the common-cut identities directly from finite words, using the
spatial first moment proved in Section~\ref{sec:geometry}.

For two starting sites \(x,y\) on the same layer, we distinguish the
independent conditioned law \(P_x^+\otimes P_y^+\) from the annealed law
\(P_{x,y}^{++}\) of two walks in the same environment, conditioned on both
no-drop events \(D\). Let \(P_{x,y}^{\rm sh}\) denote the latter pair law
before conditioning, and put
\[
p_2(x,y)=E[q(x)q(y)].
\]
The positivity of \(q\) established in Section~\ref{sec:geometry} gives a
uniform lower bound
\[
\underline p_2:=\inf_{x_1=y_1}p_2(x,y)>0.
\]
Indeed, for every \(c>0\),
\[
P\bigl(\min(q(x),q(y))<c\bigr)\le 2P(q(0)<c),
\]
so a sufficiently small \(c\) makes \(q(x)q(y)\ge c^2\) on an event of
uniformly positive probability. Product-field mixing also gives
\[
p_2(x,y)\longrightarrow p^2
\qquad\text{as }\|x-y\|\longrightarrow\infty,\quad x_1=y_1.
\]
To verify this directly, approximate \(q(0)\) in \(L^1(P)\) by a bounded
function of finitely many rows and translate the approximation to \(x\)
and \(y\). The approximating functions are independent once their
translated row sets are disjoint.

Each marginal of the shared conditioned law is dominated uniformly by
the corresponding single conditioned law:
\[
P_{x,y}^{++}(X^{(1)}\in A)
 =\frac{E[P_{x,\omega}(A\cap D)q(y)]}{p_2(x,y)}
 \le \frac{p}{\underline p_2}P_x^+(A),
\]
and likewise for the second walk. With superscripts distinguishing the
walks, define
\[
\Delta_j=Y_j^{(2)}-Y_j^{(1)},\qquad
Z_j=y_2-x_2+\Delta_j.
\]
A \emph{common true level} is a level that is true for both walks at their
respective first hits. Heights are measured from the common starting
layer, and level zero is included under either conditioned law.

\begin{lemma}[Common words and their continuation laws]
\label{fluct-common-words}
Under either conditioned pair law, common true levels occur infinitely
often. At each such level the translated pair continuation is fresh:
under independent environments it has law \(P_0^+\otimes P_0^+\), while
under a shared environment, from terminal sites \(x',y'\), it has law
\(P_{x',y'}^{++}\) before translation. The latter transition law depends
on the terminal pair of sites.

These continuation rules remain valid at finite stopping indices of the
chain of common words, when the information retained consists of the
traversed words and their path prefixes.
\end{lemma}

\begin{proof}
Work first with the raw pair laws. Test a new integer level above both
explored path pasts, pausing each walk at its first arrival there. All
departure rows revealed by these prefixes lie strictly below the
candidate level. The conditional probability that both arrivals are
true is \(p^2\) in independent environments and \(p_2\) at the arrival
sites in a shared environment. Both probabilities have a uniform
positive lower bound.

If the candidate fails, resume the two walks, for example one step of
each per round, until one drops below the candidate. This failure is
detected in finite time. Then test a level strictly above both explored
pasts. Each walk tends to positive infinite height almost surely; this
also holds for the shared raw pair by its marginals. Thus, after every
detected failure, another candidate can be reached. The uniform success
bound proves the existence of a common true level above any prescribed
height. Conditioning on the two initial no-drop events preserves this
almost-sure conclusion.

For the continuation law, specify a pair of finite words from \(x,y\)
to their first hits \(x',y'\) of a common terminal level \(L>0\).
Call this word pair admissible when both prefixes avoid the initial
drop and every intermediate positive level has a witnessed failure of
common truth: at least one prefix has dropped below that level after its
first hit. These are exactly the required prefix conditions. Indeed,
once truth is imposed at \(L\), any failure of truth at an earlier level
must already have occurred in the specified prefixes.

Let \(A_{x,y}^{\rm sh}(w_1,w_2)\) be the raw annealed prefix weight of an
admissible word pair. Its quenched weight is the product of the two path
weights in the shared environment, including any repeated use of a row.
Every departure row lies below the terminal layer. The suffixes
constrained by truth use only rows at or above that layer, so site-row
independence gives the conditional word probability
\[
P_{x,y}^{++}\bigl((w_1,w_2)\text{ is the first common word}\bigr)
 =
 A_{x,y}^{\rm sh}(w_1,w_2)
 \frac{p_2(x',y')}{p_2(x,y)}.
\]
More generally, for any event \(B\) of the suffix pair, its joint
probability with this word pair is
\[
A_{x,y}^{\rm sh}(w_1,w_2)
\frac{p_2(x',y')}{p_2(x,y)}
P_{x',y'}^{++}(B).
\]
Here \(B\) is read with the new sites as starting sites. In particular,
the factor \(p_2(x',y')/p_2(x,y)\) is part of the shared word law; it
does not cancel in general.

In independent environments the corresponding terminal and initial
conditioning probabilities are both \(p^2\), so their ratio is one.
Translation of each walk separately then makes the relative suffix law
\(P_0^+\otimes P_0^+\), independent of the word just traversed. Enumerate
the countably many admissible finite word pairs and iterate these
factorizations to obtain the stated continuation rules. Summing over a
finite stopping index gives the same rules there. The information at
such an index does not include any unseen upper environment.
\end{proof}

In independent environments, let \(L\ge1\) be the integer width of a
common word and let \(S\) be its increment of \(\Delta\). The successive
pairs \((L_i,S_i)\), with their path decorations, are iid by
Lemma~\ref{fluct-common-words}. Exchanging the walks shows that \(S\) is
symmetric. In what follows, expectations of \(L,S\) refer to this
independent common-word law.

\begin{lemma}[Common-word moments and local tightness]
\label{fluct-local-tightness}
The independent common-word variables satisfy
\[
1\le \bar L:=\mathbb E L\le p^{-2},\qquad
\mathbb E|S|<\infty,\qquad \Pr(S\ne0)>0.
\]
For a single conditioned walk, and for either marginal of a shared
conditioned pair, the increments \(Y_{h+j}-Y_h\) are tight uniformly in
the deterministic level \(h\) when \(j\) ranges over any fixed bounded
set of nonnegative integers.

If \(\sigma(H)\) is the first common true level at or above \(H\) under
\(P_{x,y}^{++}\), then
\[
\sigma(H)-H,\quad
Y^{(i)}_{\sigma(H)}-Y^{(i)}_H\ (i=1,2)
\quad\text{are tight uniformly in }H,x,y.
\tag{6}\label{eq:6}
\]
At \(\sigma(H)\) the pair suffix has the shared conditioned law from its
new sites.
\end{lemma}

\begin{proof}
The probability that a deterministic level \(h\) is true under the
single conditioned law is \(u_h\), by the exact-truth identity in
Section~\ref{sec:geometry}. Independence therefore makes the renewal
mass for common words equal to \(u_h^2\ge p^2\). The renewal-count strong
law, including its infinite-mean version obtained by truncation, implies
\(\bar L\le p^{-2}\). The lower bound follows from \(L\ge1\).

For each of the two paths, let \(N_i\) be the number of its single words
used before the first positive common boundary. Since single widths are
positive integers, \(N_i\le L\). The event \(N_i\ge j\) is determined by
that path's first \(j-1\) words and the independent boundary sequence of
the other path. It is consequently independent of the radius of its
\(j\)-th word. The finite mean single-word radius from
Section~\ref{sec:geometry} gives
\[
\mathbb E\sum_{j=1}^{N_i}R_{{\rm s},j}^{(i)}
 =\mathbb E N_i\,\mathbb E R_{\rm s}
 \le \bar L\,\mathbb E R_{\rm s}<\infty.
\]
The absolute value of \(S\) is bounded by the sum of these two total
radii. Finally, uniform ellipticity permits two short paths with
different second-coordinate displacements to level one. Imposing fresh
truth there gives positive probability of an exact common cut with
nonzero \(S\).

We record the renewal estimate that supplies local tightness. For either
the single renewal sequence or the independent common sequence, let
\(L'\) be its word width and \(A\) any almost surely finite size of the
word, such as its maximal displacement. If \(A_{\rm word}\) is the size
of the word based at the last boundary at or below a deterministic
level \(h\), then
\[
\Pr(A_{\rm word}>s)
 \le \sum_{j=0}^h\Pr(L'>j,\ A>s)
 \le \mathbb E[L'\mathbf1_{\{A>s\}}].
\]
To obtain the first inequality, sum over the possible preceding
boundary levels \(h-j\), use independence of the next word, and bound
each renewal mass by one. Since \(\mathbb E L'<\infty\), the final
quantity tends to zero as \(s\to\infty\). The same reasoning with
\(A=L'\) gives tightness of the covering width. Thus first-hit positions
at a deterministic layer differ from their preceding boundary
positions by tight amounts. A fixed interval of integer levels meets
only the words covering those levels, so it also has tight total
displacement, uniformly in its base. Shared marginals inherit this
statement from their uniform domination by the single conditioned law.

It remains to prove \eqref{eq:6}. Test the \(m\) levels
\[
h=H,H+B,\ldots,H+(m-1)B.
\]
Call failure at \(h\) \emph{visible} if, after first hitting \(h\), some
walk drops below it before its first hit of \(h+B\). In the raw law,
the probability that all \(m\) tests have visible failure is at most
\((1-\underline p_2)^m\). Indeed, at each pair of first-hit arrivals to
\(h\), all previous visible-failure tests are already determined, while
both fresh truths at \(h\) would preclude the next failure. Passing to
the conditioned pair law costs at most \(\underline p_2^{-1}\).

If \(h\) is not common true but has no visible failure, some walk drops
below \(h\) after its first hit of \(h+B\). That walk has no single true
boundary between \(h\) and \(h+B\), inclusive; its word covering \(h\)
therefore has width greater than \(B\). The covering-word estimate and
marginal domination bound this alternative by
\[
C\mathbb E[L_{\rm s}\mathbf1_{\{L_{\rm s}>B\}}].
\]
A union bound over the \(m\) tests now gives
\[
P_{x,y}^{++}\bigl(\sigma(H)>H+(m-1)B\bigr)
 \le \underline p_2^{-1}(1-\underline p_2)^m
      +Cm\mathbb E[L_{\rm s}\mathbf1_{\{L_{\rm s}>B\}}].
\]
Choose \(m\) large and then \(B\) large. This proves uniform tightness of
the height overshoot. The fixed-interval displacement bound then proves
the other assertions in \eqref{eq:6}. Finally, the first common boundary
at or above \(H\) is a stopping index of the common-word chain, so its
suffix rule follows from Lemma~\ref{fluct-common-words}.
\end{proof}

\subsection{Truncated variance and median centering}

For \(r>0\), define
\[
m(r)=\mathbb E(S^2\wedge r^2),\qquad n(r)=r^2/m(r).
\]
These quantities are positive by \(\Pr(S\ne0)>0\). The scale \(n(r)\)
uses only the finite first moment of \(S\); a finite second moment is
not assumed.

\begin{lemma}[Scale and maximal estimates]\label{fluct-scales}
The function \(n\) is continuous and nondecreasing, tends to infinity,
and is strictly increasing for all sufficiently large \(r\). Write
\(r_s\) for its inverse on large scales, so \(n(r_s)=s\). Then
\[
n(cr)\ge c^2n(r)\ (0<c\le1),\qquad
n(r)\ge r/\mathbb E|S|.
\tag{7}\label{eq:7}
\]
For iid copies \(S_j\), let
\[
M_H^S=\max_{0\le i\le H}\left|\sum_{j=1}^iS_j\right|.
\]
Uniformly in integer \(H\ge0\) and \(z>0\),
\[
\mathbb E[(M_H^S)^2\wedge z^2]\le CHm(z),\qquad
\Pr(M_H^S>z)\le CH/n(z).
\tag{8}\label{eq:8}
\]
\end{lemma}

\begin{proof}
The ratio
\(m(r)/r^2=\mathbb E[(S^2/r^2)\wedge1]\) is continuous and
nonincreasing. There is a bounded interval of strictly positive values
of \(|S|\) having positive probability; its contribution is strictly
decreasing once \(r\) exceeds that interval. This proves the
monotonicity and eventual strict monotonicity of \(n\). Monotonicity of
\(m\) proves the first inequality in \eqref{eq:7}, and
\(S^2\wedge r^2\le r|S|\) proves the second, hence \(n(r)\to\infty\).

For \eqref{eq:8}, replace \(S_j\) by
\(S_j\mathbf1_{\{|S_j|\le z\}}\). These truncated variables have mean
zero by symmetry. The probability of any replacement among the first
\(H\) variables is at most \(H\Pr(|S|>z)\). Off this event the original
partial sums agree with the truncated sums, whose expected squared
maximum is at most
\(4H\mathbb E[S^2\mathbf1_{\{|S|\le z\}}]\) by the square maximal
inequality. On the replacement event, the capped squared maximum is
at most \(z^2\). Adding these bounds gives the first inequality in
\eqref{eq:8}; Markov's inequality for the capped square gives the
second.
\end{proof}

\begin{lemma}[Exact-level comparison]\label{fluct-exact-comparison}
For every integer \(H\ge0\) and \(z>0\),
\[
P_0^+(W_H>z)\le C\Pr(M_H^S>z).
\tag{9}\label{eq:9}
\]
No moment of the maximal decoration inside a common word is needed.
\end{lemma}

\begin{proof}
In a raw first walk, select the first level \(h\le H\), if any, such
that \(T_h<T_{-1}\) and \(|Y_h-b_h|>z\). Let \(A\) be the event that
such a level exists. Conditional on this prefix and its selected
level, an independent raw second walk reaches \(h\) before dropping,
with its displacement on the opposite side of the median \(b_h\),
with probability at least \(p/2\). Here the conditioned no-drop law
only lower bounds a prefix probability; it is not imposed as a
separate test on that second walk.

At the two first-hit endpoints, impose both fresh truths. Their
probability factor is \(p^2\). The resulting event implies the two
initial no-drop events and has \(h\) as an exact common true level,
with \(|\Delta_h|>z\). Dividing its raw probability by \(p^2\), its
probability under the independent conditioned pair law is at least
\(P_0(A)p/2\). At a common cut of height at most \(H\), the difference
is a partial sum of at most \(H\) common-word increments, so
\[
\Pr(M_H^S>z)\ge \frac p2P_0(A)
 \ge \frac{p^2}{2}P_0^+(W_H>z).
\]
This proves \eqref{eq:9}. All of these adaptive prefix transfers can
also be obtained by summing the corresponding finite stopped-word
identities separately over the possible levels \(h\).
\end{proof}

\begin{lemma}[Median linearization on short blocks]
\label{fluct-median-linearization}
For sufficiently small \(t>0\), uniformly in
\(1\le H\le tn(r)\),
\[
\max_{j\le H}|b_j-jb_H/H|\le C\sqrt t\,r.
\tag{10}\label{eq:10}
\]
\end{lemma}

\begin{proof}
By \eqref{eq:7}--\eqref{eq:9}, choose a fixed large \(C_0\), then
\(t\) small enough that \(C_0\sqrt t\le1\), to obtain
\[
P_0^+(W_{2H}>C_0\sqrt t\,r)<p/4.
\]
For \(i+j\le2H\), the event that \(i\) is exactly true under \(P_0^+\)
has probability \(u_i\ge p\), and its suffix has the fresh conditioned
law. There is therefore positive probability on this event that
\(Y_i\), \(Y_{i+j}\), and \(Y_{i+j}-Y_i\) are all within
\(C_0\sqrt t\,r\) of \(b_i\), \(b_{i+j}\), and \(b_j\), respectively.
For the first two requirements use the preceding maximal bound, and
for the third use the same bound on the fresh suffix. Consequently,
\[
b_{i+j}=b_i+b_j+O(\sqrt t\,r)
\qquad(i+j\le2H),
\]
uniformly in the indices.

Set \(d_j=b_j-jb_H/H\) for \(0\le j<H\). Apply the last relation to
\(j+j\), and, when \(2j\ge H\), also to \(H+(2j-H)\). This gives
\[
2d_j=d_{2j\bmod H}+O(\sqrt t\,r)
\]
with a uniform error. Taking the maximum absolute value over
\(0\le j<H\) bounds that maximum by half itself plus
\(C\sqrt t\,r\). Since the error at \(j=H\) is zero, this proves
\eqref{eq:10}.
\end{proof}

\subsection{Change of conditioning and short-block survival}

The estimates above concern \(P_x^+\), which conditions the joint
annealed law on no drop. Under that law environments are weighted by
\(q(x)/p\). For quenched kernel estimates we also need the mixture
\[
\widetilde P_x
 =\int P_{x,\omega}(\,\cdot\mid D)\,P(d\omega),
\]
which keeps the original environment law and conditions the path in
each environment separately. We first transfer large-scale fluctuation
bounds to this mixture. For blocks of height \(H\le tn(r)\), a fresh
restart after the first loss of mass will retain positive kernel mass
outside an environment exception of probability \(O(t^2)\).

\begin{lemma}[Transfer between conditioned laws]
\label{fluct-conditioning-transfer}
Along any sequences of integers \(H\ge0\) and real numbers
\(z\to\infty\),
\[
\limsup\widetilde P_0(W_H>z)
 \le C\limsup P_0^+(W_H>z/C)
\tag{11}\label{eq:11}
\]
with a fixed constant \(C\). If also \(H\to\infty\), the same bound
holds with the left-hand side evaluated under the \(P\)-average of
the normalized successful-prefix kernel \(K_H/a_H\) at the origin.
\end{lemma}

\begin{proof}
On the joint environment-and-path space,
\[
\frac{d\widetilde P_0}{dP_0^+}=\frac{p}{q(0)}.
\]
This density is integrable but need not be bounded. For any prescribed
error, truncate it and then approximate the truncation by a bounded
nonnegative environment function \(g\) depending on finitely many
rows. Thus \(g\) approximates \(p/q(0)\) in \(L^1(P_0^+)\), and
\(E_0^+g\) tends to one as the approximation error tends to zero.
We will remove the finite initial segment that can be affected by \(g\).

First check the deterministic centering error caused by such a removal.
If \(\limsup P_0^+(W_H>z/10)\ge0.3\), Equation~\eqref{eq:11}
is immediate with a sufficiently large fixed \(C\). Otherwise these
probabilities are eventually less than \(0.3\). For every fixed
integer \(s_0\), local tightness then gives
\[
|b_{h+s}-b_h|\le0.3z
\qquad(0\le s\le s_0,\ h+s\le H)
\]
for all sufficiently large indices in the sequence. Indeed, each of
\(|Y_h-b_h|\le z/10\) and
\(|Y_{h+s}-b_{h+s}|\le z/10\) has probability greater than \(0.7\).
By Lemma~\ref{fluct-local-tightness}, the probability of
\(|Y_{h+s}-Y_h|\le z/10\) tends to one uniformly over these indices.
The three events have a nonempty intersection, which proves the
claimed deterministic inequality by the triangle inequality.

Let \(s\) be the first positive single true boundary strictly above
every row used by \(g\). Its index in the single-word chain is a finite
stopping index. For a sufficiently accurate approximation, \(E_0^+g>0\).
Under the probability measure with density \(g/(E_0^+g)\) relative to
\(P_0^+\), the translated suffix at \(s\) has law \(P_0^+\),
independently of the prefix. To verify this, apply exact-word
factorization separately to each finite prefix through the selected
boundary. Both its path weight and \(g\) use only rows below that
boundary; all rows for the conditioned suffix are fresh. This remains
true for rows used by \(g\) that the prefix has not visited. The
separation is by height, not by which rows the walk happened to reveal.

For fixed \(g\), first choose a deterministic \(s_0\) so large that
\(s\le s_0\) except for an arbitrarily small \(g\)-weighted error.
As \(z\to\infty\), we may further restrict to
\[
\max_{j\le s}|Y_j|\le z/10,
\qquad \max_{j\le s_0}|b_j|\le z/10,
\]
with a further error tending to zero. On these restrictions no offense
\(|Y_j-b_j|>z\) can occur before \(s\). After \(s\), write the
relative suffix displacement as \(Y'_u=Y_{s+u}-Y_s\). The centering
estimate above gives, for \(s<j\le H\),
\[
|Y_j-b_j|
\le |Y_s|+|Y'_{j-s}-b_{j-s}|+|b_{j-s}-b_j|
\le0.4z+|Y'_{j-s}-b_{j-s}|.
\]
Thus \(W_H>z\) requires the fresh suffix to have its own
\(W_H>z/10\). If \(H\le s\), no offense remains at all.

The \(g\)-weighted probability is consequently bounded in the limsup
by \(E_0^+g\) times \(\limsup P_0^+(W_H>z/10)\), plus the
restriction errors. Return to \(\widetilde P_0\) using the
\(L^1\) density approximation, and send all approximation and
restriction errors to zero. Enlarging the same fixed \(C\) to include
the earlier trivial case proves Equation~\eqref{eq:11}.

Finally, in a fixed environment the event \(D\) is a subevent of
\(\{T_H<T_{-1}\}\). On first-hit prefixes, changing between these
two conditional laws costs total variation at most \(1-q(0)/a_H\).
Therefore the distance between \(\widetilde P_0\) on those prefixes
and the \(P\)-average of \(K_H/a_H\) is at most
\[
E(1-q(0)/a_H)\longrightarrow0,
\]
since \(a_H\downarrow q(0)>0\) almost surely. This proves the last
assertion.
\end{proof}

\begin{proposition}[Short-block survival]
\label{fluct-short-block}
For every \(v_0>0\) there are constants \(C_{v_0}<\infty\) and
\(t_{v_0}>0\) such that, for each fixed \(0<t\le t_{v_0}\), some
\(\delta>0\) satisfies, for all sufficiently large \(r\), every
integer \(1\le H\le tn(r)\), and every probability measure \(\pi\)
on a starting layer,
\[
P\left(
 \pi K_H\{W_H\le v_0r,\
       \max_{j\le H}|Y_j-jb_H/H|\le v_0r\}<\delta
 \right)
 \le C_{v_0}t^2.
\tag{12}\label{eq:12}
\]
The environment is fresh with product law, and path displacements
are measured from each path's own starting site.
\end{proposition}

\begin{proof}
Fix \(v_0>0\), set \(w_0=v_0/8\), and write
\(\theta=b_H/H\). The proof uses one estimate for loss of mass in a
narrower slope tube, then restarts that estimate after its first failure.
All stopping tests will use finite-height kernels.

\paragraph{One possible loss of mass.}
Choose \(t_{v_0}\) small enough that Equation~\eqref{eq:10} gives,
for sufficiently large \(r\), uniformly over the stated \(H\),
\[
\max_{j\le H}|b_j-j\theta|\le w_0r/2,
\qquad |\theta|\le w_0r.
\]
The second inequality follows from the first at \(j=1\), since
\(b_1\) is fixed. Equations~\eqref{eq:7}--\eqref{eq:9} and
\eqref{eq:11} also give
\[
E[P_{0,\omega}(W_H>w_0r/2\mid D)]\le C_{v_0}t
\]
uniformly in \(1\le H\le tn(r)\) for large \(r\). Indeed, a
violating sequence of \(H,r\) would contradict the limsup transfer
in Equation~\eqref{eq:11}, whose right side is bounded by a constant
times \(H/n(r)\). The scale inequality in Equation~\eqref{eq:7}
absorbs the fixed change in threshold.

Choose \(\delta_0\in(0,1/4)\) so that
\(P(q(0)<4\delta_0)\le t\). Call a starting site \(x\) favorable
when
\[
q(x)\ge4\delta_0,
\qquad P_{x,\omega}(W_H\le w_0r/2\mid D)\ge1/2.
\]
The expected \(\pi\)-fraction of unfavorable sites is at most
\(C_{v_0}t\). Markov's inequality shows that, except on an environment
event of probability at most \(C_{v_0}t\), favorable sites carry at
least half the initial mass. From each favorable site the raw mass of
paths on \(D\) in the median tube is at least \(2\delta_0\).
Their first-hit prefixes lie in the slope tube of width \(w_0r\) at
every shorter horizon. Hence, except on that one environment event,
\[
M_u^\pi:=
\pi K_u\left\{\max_{j\le u}|Y_j-j\theta|\le w_0r\right\}
\ge\delta_0
\qquad(0\le u\le H).
\]
In particular the probability that any of these inequalities fails is
at most \(C_{v_0}t\), uniformly in the initial probability \(\pi\).
No independence among the environments seen from different starting
sites was used.

\paragraph{A fresh restart after the first failure.}
Test the masses \(M_u^\pi\) in increasing order, and let \(s\) be
the first level at which the displayed inequality fails. There is no
failure at level zero, since \(M_0^\pi=1\). Each test is measurable
from rows strictly below its tested layer, so \(s\le H\), when finite,
is an environment stopping layer.

The passing submeasure at layer \(s-1\) has mass at least
\(\delta_0\). Extend every one of its prefixes by one direct
\(+e_1\) step. Uniform ellipticity leaves mass at least
\(\kappa\delta_0\). The extension has slope error at most
\(2w_0r\), since the earlier error is at most \(w_0r\) and
\(|\theta|\le w_0r\). Both detection of the failure and this bridge
use only rows below layer \(s\).

Normalize the bridged endpoint submeasure to a probability \(\pi_s\)
and restart at layer \(s\). Conditional on the exposed lower rows,
\(\pi_s\) is fixed and the upper field has its original product law.
The one-failure estimate therefore applies again. Use the same slope
\(\theta=b_H/H\) and its original horizon \(H\); its simultaneous
conclusions imply the required ones through the shorter remaining
horizon \(H-s\). Replacing \(H\) by \(H-s\) in the slope would
not be justified and is unnecessary.

Thus, conditional on the first failure, the probability of a second
failure before the remaining horizon is at most \(C_{v_0}t\).
The probability of two failures is at most \(C_{v_0}t^2\), after
enlarging the constant. This multiplication uses fresh upper rows at
the stopping layer, not independence of two events in the original slab.

\paragraph{Retained mass and its tube.}
If no first failure occurs, at least \(\delta_0\) of the original
kernel mass reaches the top in the slope tube. If there is a first
failure but no second one, concatenate the bridged mass and the restarted
successful prefixes. Retain each prefix's original starting site when
doing so; restarted displacements are measured from the new endpoint.
The total slope error is at most \(2w_0r+w_0r=3w_0r\), and the
retained mass is at least \(\kappa\delta_0^2\). Local no-drop
successes are also successes above the original floor, so the
concatenated kernels are submeasures of the original \(K_H\).

Adding the deterministic median-line error \(w_0r/2\) puts these
prefixes in the required median tube as well. Both tube widths are
less than \(v_0r\). Take
\(\delta=\min\{\delta_0,\kappa\delta_0^2\}\) to obtain
Equation~\eqref{eq:12}. The infinite no-drop event served only to bound
failure probabilities; every test, bridge, and restart used the
finite-height prefix kernels.
\end{proof}

\subsection{Endpoint spread}

\begin{lemma}[Spread at the truncated-variance scale]
\label{fluct-endpoint-spread}
Under the independent conditioned pair law, there is a constant
\(c>0\) such that
\[
\begin{split}
&\Pr(|\Delta_{\lfloor n(r)\rfloor}|>cr)\ge c
       \qquad(r\ \text{large}),\\
&\liminf\Pr(|\Delta_{\lfloor tn(r)\rfloor}|>lr/4)
       \ge c_{l,\epsilon}t
 \quad\text{if }r\to\infty,\qquad
       n(r)\Pr(|S|>lr)\ge\epsilon.
\end{split}
\tag{13}\label{eq:13}
\]
In the second assertion \(0<l<1\) and \(\epsilon>0\) are fixed,
and the bound holds for each sufficiently small fixed \(t>0\),
depending on \(l,\epsilon\). Relative to any deterministic endpoint
center, a single \(P^+\)-walk consequently has a comparable
probability of a deviation, at half the corresponding separation
threshold, in at least one of the two signs.
\end{lemma}

\begin{proof}
Fix \(s>0\), put \(H=\lfloor sn(r)\rfloor\), and let
\(k=\lfloor H/\bar L\rfloor\). First we justify the approximation
\[
\Delta_H-\sum_{i=1}^kS_i=o_{\Pr}(r).
\]
By Lemma~\ref{fluct-local-tightness}, the displacement from the
preceding common boundary to the first-hit position at \(H\) is
tight. The renewal strong law puts the number of completed common
words within \(o_{\Pr}(n(r))\) of \(k\). No independence between
this count and the increments is required: on the event that its
distance from \(k\) is at most \(an(r)\), bound the sum error by
the maximal sums in the deterministic forward and backward
windows of that length around \(k\). For every fixed
\(\rho>0\), \eqref{eq:7}--\eqref{eq:8} bound the probability of
such an error exceeding \(\rho r\) by
\(O_\rho(a)+o(1)\). Let \(a\downarrow0\), and include the tight
decoration error, to obtain the approximation.

For \(s=1\), set
\[
V=\sum_{i=1}^k(S_i^2\wedge r^2).
\]
Its mean is comparable to \(r^2\), since
\(k\sim n(r)/\bar L\), and
\[
\mathbb E V^2
 \le kr^2m(r)+k^2m(r)^2=O(r^4).
\]
The second-moment inequality therefore makes \(V\) exceed a fixed
positive multiple of \(r^2\) with uniformly positive probability.
Conditional on the magnitudes \((|S_i|)\), the nonzero signs are
independent fair signs. The sum \(\sum_{i=1}^kS_i\) has conditional
variance \(Q=\sum_{i=1}^kS_i^2\ge V\) and conditional fourth
moment at most \(3Q^2\). A further second-moment inequality gives
a uniformly positive conditional probability of a displacement
of order \(\sqrt Q\), hence of order \(r\) on the preceding event.
The approximation, with a smaller displacement constant, proves
the first line of \eqref{eq:13}.

For the second line, take \(s=t\) and set
\(p_r=\Pr(|S|>lr)\). The assumptions and the definition of \(m\)
give
\[
\epsilon\le n(r)p_r\le l^{-2}.
\]
Among \(k\sim tn(r)/\bar L\) summands, the probability that exactly
one has magnitude greater than \(lr\) is therefore at least
\(c_{l,\epsilon}t\) for each sufficiently small fixed \(t>0\) and
large \(r\). Conditional on the index of that exceptional
summand, the others are independent symmetric variables
conditioned to have magnitude at most \(lr\). Their sum has
variance at most
\[
\frac{k\,\mathbb E[S^2\mathbf1_{\{|S|\le lr\}}]}{1-p_r}
 \le Ctr^2.
\]
For \(t\) sufficiently small, Chebyshev's inequality gives
probability at least \(1/2\) that their sum has magnitude at most
\(lr/2\). On this event they cannot cancel more than half the
exceptional summand. Thus the deterministic sum has magnitude
greater than \(lr/2\) with probability at least a constant times
\(t\). Keep \(t\) fixed while \(r\to\infty\); the
\(o_{\Pr}(r)\) approximation then proves the claimed bound at
threshold \(lr/4\).

Finally, if two iid endpoint displacements differ by more than
\(a\), at least one differs from any given deterministic center
by more than \(a/2\). The probability of this single-walk
absolute deviation is at least half the two-walk separation
probability. One of its two signs has at least half that mass.
This proves the final assertion after adjusting constants.
\end{proof}
